% GENERATED by `bench report` — DO NOT EDIT BY HAND; regenerate instead.
% source: /tmp/claude-0/-workspace-retrieve/7c1995ca-a1a3-4b78-9907-2c0de7f4d6fc/scratchpad/all   generated: 2026-09-15T20:57:27Z
% records: 24 {'ok': 24}   cells flagged unstable: 21
% code_version: 0fe440d4bc9013097737e225af59b71ccec94a6d   commit: 71430d7,afc2ab9   branch: dev/c4-gate-rerun,dev/c5-harness-split
% schema_version: 1,2   gpu: NVIDIA A100-SXM4-80GB   run window: 2026-09-15T08:50:50+00:00 .. 2026-09-15T13:54:29+00:00
% PROVENANCE: *** NOT CITABLE *** (CLAUDE.md rule 2). Reasons:
%   - no --gate given: these records come from no green roadmap gate
%   - record(s) produced on dev/c4-gate-rerun, dev/c5-harness-split — CLAUDE.md rule 2: harness numbers from a branch are not paper material
%   Built from pre-campaign records: roadmap D1 has not produced these.
% Для \section{Методология замеров скорости и памяти} (docs/thesis/main.tex).
\begin{itemize}
    \item \textbf{Стенд.} NVIDIA A100-SXM4-80GB, драйвер 580.159.04, CUDA 12.8, PyTorch 2.10.0+cu128, Triton 3.6.0, Python 3.11.15. Частоты SM \textbf{нельзя зафиксировать} в этом контейнере (\texttt{nvidia-smi -lgc} запрещён), поэтому каждая запись хранит выборку частоты под нагрузкой (\texttt{perf[].sm\_mhz}), и все сравнения задержек делаются при одинаковом оценщике частоты.
    \item \textbf{Изоляция.} Один процесс на $(\mathrm{dataset}, \mathrm{dim}, \mathrm{algo}, \mathrm{backend})$: ни кэш \texttt{torch.compile}, ни пул CUDA-графов, ни аллокатор не переживают смену бэкенда.
    \item \textbf{Прогрев и замер.} 50 прогревочных вызовов, затем 3 окна по $N = \mathrm{clamp}(2\,\text{с} / \tilde t, 1000, 5000)$ вызовов; каждый вызов обрамлён CUDA-событиями, окно — одним \texttt{synchronize} в конце. Запросы берутся из фиксированного пула в 4096 батчей по кругу, без сброса L2. Отчёт ведётся по окну с медианным медианом: median / mean / p95 / p99 / min / IQR, QPS (closed loop), \texttt{host\_gap}, разброс окон $(\max-\min)/\mathrm{med}$; разброс выше 5\,\% помечает ячейку \texttt{unstable}.
    \item \textbf{Режимы.} eager, graph: \texttt{eager} — число, сопоставимое с обеими статьями, \texttt{graph} (\texttt{torch.compile(mode="reduce-overhead", dynamic=False, fullgraph=True)}) — приводится рядом, никогда вместо. Захват графа проверяется: \texttt{cudagraph\_skips} обязан быть 0, иначе ячейка — ошибка, а не число.
    \item \textbf{Качество.} Один проход в режиме eager при $k_{\max} = \max(K)$, чанками по 16 запросов; метрики recall / ndcg / precision / mrr при $K \in \{100, 500, 1000\}$ накапливаются на устройстве. Эталон: точный фильтрованный fp32-оракул на фильтрующих ячейках и отложенные взаимодействия всегда.
    \item \textbf{Память.} \texttt{index\_mib} — сумма всех буферов модуля поиска вместе с подмодулем фильтра; \texttt{peak\_fwd\_mib} — \texttt{max\_memory\_allocated} за первое eager-окно.
    \item \textbf{Повторы.} Размеры батча $B \in \{1, 8, 16\}$, сиды $\{0, 1, 2\}$; по сидам приводится медиана с усами min--max.
\end{itemize}
